\clearpage\pagestyle{matSimple}\phantomsection\label{guide}\pdfbookmark[0]{Guide de l’adulte}{bm-guide}
\matPageTitre{Guide de l’adulte}{Objectifs, manipulations, mots à faire dire et réponses}
{\fontsize{9.5}{13.5}\selectfont Pour chaque activité : l’objectif tiré du programme, ce que l’on peut faire avec de vrais objets avant la fiche, les mots que l’enfant est invité à dire, puis la réponse attendue. Les images en vert montrent une réponse correcte ; d’autres réponses peuvent être justes lorsque la consigne le permet.\par}
\par\needspace{62mm}\vspace{6mm}{\otSemi\fontsize{9}{11}\selectfont\color{matAccent}\otCaps[8]{Partie 1 · Plus, moins, autant}}\par\vspace{1mm}
\guideEntree{act-ps-1}{1}{Beaucoup ou peu ?}{Comparer globalement deux quantités très différentes et dire « plus que », « moins que ».}{Mettre huit jetons dans un pot transparent et deux dans un autre ; demander, sans compter, où il y en a le plus.}{« Ici, il y a beaucoup de pommes. Là, il y en a peu : il y en a moins. »}{Le panier de droite : il contient huit pommes, celui de gauche deux seulement. On ne demande pas de compter.}{ps-01}
\guideEntree{act-ps-2}{2}{Le moins d’oiseaux}{Comparer globalement deux quantités très différentes et dire « moins que ».}{Poser beaucoup de cubes sur une assiette et seulement quelques-uns sur une autre ; demander où il y en a moins.}{« Sur cet arbre, il y a moins d’oiseaux. Sur l’autre, il y en a plus. »}{L’arbre de droite (trois oiseaux contre sept). La comparaison se fait d’un coup d’œil, sans compter.}{ps-02}
\guideEntree{act-ps-3}{3}{Une carotte pour chaque lapin}{Comparer deux collections par correspondance terme à terme.}{Mettre la table pour des peluches : une assiette devant chaque peluche, puis vérifier qu’aucune n’est oubliée.}{« Chaque lapin a sa carotte. Il y a autant de lapins que de carottes. »}{Oui : chaque lapin reçoit une carotte et aucune carotte ne reste seule. N’importe quel appariement convient.}{ps-03}
\guideEntree{act-ps-4}{4}{Il manque un gâteau}{Comparer par correspondance terme à terme, même avec des collections plus grandes que les nombres connus.}{Distribuer un gâteau (un jeton) par assiette et constater qu’une assiette reste vide.}{« Il y a une assiette sans gâteau. Il manque un gâteau. »}{Il y a cinq assiettes et quatre gâteaux : une assiette reste vide, l’enfant dessine un gâteau pour elle.}{ps-04}
\guideEntree{act-ps-5}{5}{Une chaise pour chaque ours}{Comparer deux collections par correspondance terme à terme.}{Jeu des chaises avec des peluches : un ours par chaise, puis repérer celui qui reste debout.}{« Il y a plus d’ours que de chaises. Un ours n’a pas de chaise. »}{Quatre ours et trois chaises : un ours reste sans chaise. Celui que l’on entoure dépend des liens tracés.}{ps-05}
\guideEntree{act-ps-6}{6}{Éléphants et fourmis}{Comprendre qu’une quantité ne dépend ni de la taille ni de la place des objets.}{Aligner trois gros objets et trois petits objets serrés ; demander s’il y en a autant, puis vérifier en les appariant.}{« Les éléphants sont gros, les fourmis sont petites, mais il y en a autant : trois et trois. »}{Oui, il y en a autant : trois éléphants et trois fourmis. La taille des animaux ne change pas la quantité.}{ps-06}
\par\needspace{62mm}\vspace{6mm}{\otSemi\fontsize{9}{11}\selectfont\color{matAccent}\otCaps[8]{Partie 2 · Un, deux, trois}}\par\vspace{1mm}
\guideEntree{act-ps-7}{7}{Un seul}{Percevoir globalement une petite quantité : un, deux.}{Montrer des boîtes contenant un ou deux objets ; demander de montrer celles où il y a « un seul » objet.}{« Là, il y a un ballon : un seul. Là, il y a deux poissons. »}{Les 3 boîtes du ballon, de la voiture et de l’étoile.}{ps-07}
\guideEntree{act-ps-8}{8}{Deux}{Reconnaître des collections de deux objets de tailles et de dispositions différentes.}{Faire des paires d’objets très différents (deux cuillères, deux voitures, une grosse et une petite balle).}{« Il y a deux coccinelles : une grosse et une petite. Deux, c’est un et encore un. »}{Les boîtes des canards, des coccinelles et des clés : trois boîtes. La grosse et la petite coccinelle font bien deux.}{ps-08}
\guideEntree{act-ps-9}{9}{Trois}{Reconnaître des collections de trois objets, quelle que soit leur disposition.}{Montrer trois jetons alignés, puis en triangle, puis en vrac : c’est toujours trois.}{« Un, et encore un, et encore un : trois. »}{Les boîtes des poissons, des champignons et des ballons. Les ballons ont des tailles différentes mais sont bien trois.}{ps-09}
\guideEntree{act-ps-10}{10}{Un de plus}{Réaliser une collection contenant un objet de plus qu’une collection donnée.}{Donner une voiture, puis dire : « Je me suis trompé, j’en voulais deux » ; l’enfant en ajoute une.}{« Un, et encore un, ça fait deux. Deux, et encore un, ça fait trois. »}{Deux ronds, puis trois ronds, puis quatre ronds.}{ps-10}
\guideEntree{act-ps-11}{11}{Mes doigts}{Associer une quantité à sa représentation avec les doigts.}{Montrer deux, puis trois doigts ; l’enfant prend autant de jetons, puis montre lui-même ses doigts.}{« Deux doigts, deux pommes. Trois doigts, trois pommes. »}{Deux doigts avec la boîte de deux pommes, trois doigts avec celle de trois, un doigt avec celle d’une pomme.}{ps-11}
\guideEntree{act-ps-12}{12}{Le dé et les bonbons}{Associer une constellation de points à la quantité correspondante.}{Lancer un grand dé et prendre autant de jetons qu’il y a de points.}{« Trois points sur le dé, trois bonbons dans la boîte. »}{Le dé à trois points avec la boîte de trois bonbons, celui à un point avec la boîte d’un bonbon, celui à deux points avec celle de deux.}{ps-12}
\guideEntree{act-ps-13}{13}{Autant que sur la carte}{Constituer une collection de même quantité qu’une représentation de points.}{Situation des boîtes : poser dans chaque boîte autant de jetons que de points sur son étiquette.}{« Sur la carte, il y a deux points ; je dessine deux ronds. »}{Deux ronds, puis un rond, puis trois ronds.}{ps-13}
\guideEntree{act-ps-14}{14}{Les doigts des deux mains}{Composer deux et trois avec les doigts des deux mains.}{Lever un doigt sur chaque main et demander combien cela fait ; recommencer avec deux et un.}{« Un et un, ça fait deux. Deux et un, ça fait trois. Un et deux, ça fait trois aussi. »}{Deux, puis trois, puis trois : l’ordre des mains ne change pas le total.}{ps-14}
\guideEntree{act-ps-15}{15}{Des trains de deux couleurs}{Décomposer deux et trois : deux, c’est un et un ; trois, c’est un et deux.}{Faire des tours de deux ou trois cubes de deux couleurs et comparer les différentes tours obtenues.}{« Mon train a un wagon rouge et deux wagons bleus : trois, c’est un et deux. »}{Deux wagons : un rouge et un bleu. Trois wagons : un et deux, ou deux et un. Toutes les solutions qui utilisent les deux couleurs conviennent.}{ps-15}
\guideEntree{act-ps-16}{16}{Quatre}{Dénombrer une collection jusqu’à quatre, en déplaçant ou en pointant les objets.}{À aborder seulement quand deux et trois sont bien compris : compter des jetons en les déplaçant un à un.}{« Un, deux, trois, quatre : il y en a quatre. »}{Les boîtes des fraises, des jetons violets et des étoiles.}{ps-16}
\guideEntree{act-ps-17}{17}{Les chiffres 1, 2, 3}{Associer une écriture chiffrée à une quantité (jusqu’à trois).}{Montrer le chiffre écrit sur une carte, le nommer, et le faire correspondre à des collections de jetons.}{« Ça, c’est le chiffre deux : il faut deux étoiles. »}{2 avec la boîte de deux étoiles, 1 avec celle d’une étoile, 3 avec celle de trois étoiles.}{ps-17}
\guideEntree{act-ps-18}{18}{La chenille qui compte}{Réciter la comptine numérique de un à six, de façon ordonnée, en partant de un.}{Réciter la comptine en montant un escalier, une marche par mot-nombre, puis en tapant dans les mains.}{« Un, deux, trois, quatre, cinq, six. » Un mot pour chaque rond, sans en sauter.}{Six ronds : l’enfant dit un mot-nombre par rond, de un à six. La récitation est orale ; on ne fait pas écrire les nombres.}{ps-18}
\par\needspace{62mm}\vspace{6mm}{\otSemi\fontsize{9}{11}\selectfont\color{matAccent}\otCaps[8]{Partie 3 · Petits problèmes}}\par\vspace{1mm}
\guideEntree{act-ps-19}{19}{Les ours dans la boîte}{Résoudre un petit problème d’ajout en se représentant l’action.}{Jouer la scène avec une vraie boîte et des peluches : en mettre deux, en ajouter une, fermer la boîte et demander combien il y en a.}{« Il y avait deux ours, j’en mets encore un : il y en a trois. »}{Trois ours : le dé à trois points.}{ps-19}
\guideEntree{act-ps-20}{20}{La boîte de crayons}{Résoudre un petit problème de retrait en se représentant l’action.}{Mettre quatre crayons dans une boîte opaque, en retirer deux devant l’enfant, fermer et demander combien il en reste ; ouvrir pour vérifier.}{« Il y en avait quatre, j’en ai enlevé deux : il en reste deux. »}{Il reste deux crayons : la main qui lève deux doigts.}{ps-20}
\guideEntree{act-ps-21}{21}{Les oiseaux s’envolent}{Résoudre un petit problème de retrait en percevant la solution.}{Jouer la scène avec trois peluches sur un banc : l’une s’en va, combien en reste-t-il ?}{« Il y avait trois oiseaux. Un est parti. Il en reste deux. »}{Il reste deux oiseaux : la carte à deux points.}{ps-21}
\guideEntree{act-ps-22}{22}{Une coccinelle arrive}{Résoudre un petit problème d’ajout : un et encore un.}{Poser un jeton sur une feuille de papier, en ajouter un autre et demander combien il y en a.}{« Une coccinelle, et encore une : deux coccinelles. »}{Deux coccinelles : la main qui lève deux doigts.}{ps-22}
\par\needspace{62mm}\vspace{6mm}{\otSemi\fontsize{9}{11}\selectfont\color{matAccent}\otCaps[8]{Partie 4 · Les formes}}\par\vspace{1mm}
\guideEntree{act-ps-23}{23}{Pareil que le modèle}{Reconnaître des objets de même forme, quelles que soient leur couleur et leur taille.}{Toucher et regarder des formes en bois ; chercher dans la classe des objets qui ont la même forme qu’un disque.}{« Celui-là est rond comme le modèle, même s’il est rouge et petit. » L’adulte dit « disque ».}{Les quatre disques, quelle que soit leur couleur ou leur taille.}{ps-23}
\guideEntree{act-ps-24}{24}{Ranger les formes}{Trier et classer des formes planes selon leur forme, malgré la couleur, la taille et l’orientation.}{Ranger des formes en bois ou en carton dans trois boîtes marquées d’un disque, d’un carré et d’un triangle.}{« Ce carré va avec les carrés, même s’il est tourné. »}{Les deux disques dans la boîte du disque, les deux carrés dans celle du carré, les deux triangles dans celle du triangle.}{ps-24}
\guideEntree{act-ps-25}{25}{Chaque pièce dans son trou}{Percevoir qu’une forme reste la même quand on la tourne ; encastrer des objets.}{Jouer avec un vrai jeu d’encastrement : tourner la pièce jusqu’à ce qu’elle entre.}{« Si je tourne le cœur, il entre dans le trou du cœur. »}{Le cœur va dans le trou en forme de cœur, le triangle dans le trou triangulaire, l’étoile dans le trou en étoile, même si les pièces sont tournées.}{ps-25}
\guideEntree{act-ps-26}{26}{Les deux moitiés}{Reconstituer une image coupée en deux (assemblage de deux pièces).}{Couper en deux des images de catalogue et les faire réassembler.}{« Cette moitié de pomme va avec l’autre moitié de pomme. »}{La pomme avec la pomme, le poisson avec le poisson, la maison avec la maison.}{ps-26}
\guideEntree{act-ps-27}{27}{La pièce qui manque}{Compléter un assemblage de quatre pièces à partir d’un modèle.}{Faire un puzzle de quatre grosses pièces ; en cacher une et la chercher parmi d’autres.}{« Il manque le bas à droite : c’est la pièce avec la bouche de l’ours. »}{La pièce du milieu : le bas à droite de l’ours. La première est le haut de l’ours, la troisième vient du lapin.}{ps-27}
\guideEntree{act-ps-28}{28}{Le carrelage}{Reproduire un assemblage de quatre formes à partir d’un modèle (à une autre taille).}{Reproduire avec quatre carrés de couleur un petit modèle posé sur la table.}{« En haut, rouge puis bleu. En bas, jaune puis vert. »}{Rouge en haut à gauche, bleu en haut à droite, jaune en bas à gauche, vert en bas à droite.}{ps-28}
\par\needspace{62mm}\vspace{6mm}{\otSemi\fontsize{9}{11}\selectfont\color{matAccent}\otCaps[8]{Partie 5 · Long et court}}\par\vspace{1mm}
\guideEntree{act-ps-29}{29}{Le serpent le plus long}{Comparer visuellement deux longueurs très différentes.}{Comparer deux cordes ou deux bandes de papier très différentes, puis deux bâtons.}{« Ce serpent est plus long. L’autre est plus court. »}{Le serpent vert, en haut.}{ps-29}
\guideEntree{act-ps-30}{30}{Le crayon le plus long}{Comparer des longueurs quand les objets partent du même point.}{Poser trois crayons contre un bord de table pour qu’ils partent du même endroit, et comparer leurs bouts.}{« Ils partent tous du même trait. Celui-ci va le plus loin : c’est le plus long. »}{Le crayon du milieu.}{ps-30}
\guideEntree{act-ps-31}{31}{Même longueur}{Reconnaître un objet de même longueur qu’un objet donné.}{Chercher, parmi des bandes de papier, celle qui a la même longueur qu’une bande donnée, en les mettant côte à côte.}{« Ce ruban est aussi long que celui-là : ils ont la même longueur. »}{Le petit rouge avec le petit orange, le grand bleu avec le grand violet, le moyen jaune avec le moyen vert.}{ps-31}
\guideEntree{act-ps-32}{32}{La plus longue, la plus courte}{Comparer des longueurs ; préparer le rangement du plus long au plus court.}{Construire un escalier avec trois briques de longueurs différentes, de la plus longue à la plus courte.}{« Celle-ci est la plus longue. Celle-là est la plus courte. »}{La brique du milieu de la page est la plus longue (rouge) ; celle du haut est la plus courte (bleue).}{ps-32}
\par\needspace{62mm}\vspace{6mm}{\otSemi\fontsize{9}{11}\selectfont\color{matAccent}\otCaps[8]{Partie 6 · Les motifs}}\par\vspace{1mm}
\guideEntree{act-ps-33}{33}{Le collier}{Compléter un motif répétitif très simple (alternance de deux couleurs).}{Enfiler des perles en alternant deux couleurs, puis retirer une perle et demander laquelle il faut remettre.}{« Rouge, bleu, rouge, bleu… Ici, il manque une bleue. »}{Dans l’ordre : la quatrième perle est bleue, la septième rouge, la neuvième rouge.}{ps-33}
\guideEntree{act-ps-34}{34}{La même rangée}{Reproduire à l’identique un motif répétitif.}{Recopier une rangée de cubes de deux couleurs posée sur la table, cube par cube.}{« Rouge, jaune, rouge, jaune… C’est toujours pareil. »}{Rouge, jaune, rouge, jaune, rouge, jaune.}{ps-34}
\guideEntree{act-ps-35}{35}{Ce qui est caché}{Compléter un motif répétitif : anticiper l’élément caché.}{Aligner étoile, disque, étoile, disque avec des cartes ; en retourner une et la faire deviner.}{« Étoile, rond, étoile, rond… Sous la carte, il y a une étoile. »}{Une étoile : après le rond vient toujours une étoile.}{ps-35}
\guideEntree{act-ps-36}{36}{Le train à copier}{Reproduire à l’identique un motif répétitif.}{Copier un train de cubes de deux couleurs, d’abord en le regardant, puis de mémoire.}{« Vert, jaune, vert, jaune, vert. »}{Vert, jaune, vert, jaune, vert, derrière la locomotive.}{ps-36}
\par\needspace{62mm}\vspace{6mm}{\otSemi\fontsize{9}{11}\selectfont\color{matAccent}\otCaps[8]{Partie 7 · Se repérer}}\par\vspace{1mm}
\guideEntree{act-ps-37}{37}{Dans le bocal}{Comprendre et utiliser les mots « dans » et « dehors ».}{Jouer à mettre des objets dans une boîte ou dehors, en le disant à voix haute.}{« Ce bonbon est dans le bocal. Celui-là est dehors, sur la table. »}{Les trois bonbons du bocal. Les deux autres sont dehors, sur la table.}{ps-37}
\guideEntree{act-ps-38}{38}{Sur ou sous la table}{Comprendre et utiliser les mots « sur » et « sous ».}{Poser une peluche sur la table, puis sous la table, et le dire chaque fois.}{« Ici, le chat est sur la table. Là, il est sous la table. »}{Le chat de droite, sous la table.}{ps-38}
\guideEntree{act-ps-39}{39}{Devant ou derrière}{Comprendre et utiliser les mots « devant » et « derrière ».}{Se cacher derrière une chaise, puis se mettre devant ; faire dire où l’on est.}{« Ce lapin est derrière l’arbre : on ne le voit pas en entier. L’autre est devant. »}{Le lapin de gauche, à moitié caché par l’arbre.}{ps-39}
\guideEntree{act-ps-40}{40}{Le chemin de la souris}{Suivre et reproduire un parcours simple.}{Faire un parcours dans la cour ou dans la salle de jeux, puis le refaire avec une petite voiture sur une feuille.}{« Je vais tout droit, je tourne, je descends… »}{Le trait suit la route grise sans en sortir, de la souris jusqu’au fromage.}{ps-40}
